\section{Objectives and Contributions}
\label{sec:introduction:objectives}
Starting from BayesFlow (in Section \ref{sec:condnf}) \citep{radev2020bayesflow}, we endeavor to explore the potential of normalizing flows for system parameter identification, dynamics modeling, state estimation, and also optimal control. 
The primary objectives of this dissertation are as follows:
\begin{itemize}
    \item To investigate the strength and limitation of BayesFlow in parameter identification;
    \item To extend conditional normalizing flows for parameter identification to estimate sequential system states;
    \item To explore the capabilities of normalizing flows as control policies.
\end{itemize}
The key contributions of this dissertation can be summarized as:
\begin{itemize}
\item A comprehensive examination of the original BayesFlow framework in real-world applications, assessing its performance and generalizability.
\item Integration of BayesFlow with sensitivity analysis to enable parameter estimation across varying levels of identifiability.
\item Development of \textit{Stochastic BayesFlow}, which combines stochastic normalizing flows with BayesFlow to achieve robust parameter identification under uncertainty.
\item Extension of conditional normalizing flows through \textit{IMFlow}, introducing the concept of using flows as control policies.
\item Proposal of \textit{FlowSSM}, a novel method that leverages normalizing flows to model uncertainty in stochastic dynamical systems.
\end{itemize}